\section{Network, filesystem and industrial constraints}
\subsection{Network controls and MCP trust}
\texttt{approvedUrl()} requires an exact origin from \texttt{MCP\_ALLOWED\_ORIGINS} or the RAG REST allowlist; credential URLs, query/fragments and redirects are rejected. \texttt{guardedFetch()} validates DNS answers during the actual socket lookup and blocks configured metadata, link-local and other unsafe ranges. Approved local private addresses are deliberately usable for Docker services. This is an origin policy with address checks, not a general network sandbox.

Read-only classification currently consumes the remote server's \texttt{readOnlyHint}. Operator trust and server-side restricted credentials are essential: a malicious server can lie. MCP explicitly treats annotations as hints from a trust boundary~\cite{mcpannotations}. Core cannot prove that an arbitrary allowed ``read'' server has no side effects. ZSK's read credential independently refuses admin operations; this is stronger than relying on a name or annotation alone.
\subsection{Structured local execution}
Local task actions are disabled unless \texttt{TASK\_WORKSPACE\_ROOT} is set. \texttt{workspacePath()} resolves real paths, excludes root/outside paths, rejects protected secret/configuration directories and restricts symlinks. Workspace operations read/list/exists/write/delete have bounded content; delete requires exact approval. Writes use a temporary exclusive no-follow file, fsync and rename.

Structured commands admit only exact argument grammars for pwd, ls, cat, literal grep, bounded find and restricted git status/diff. Shell metacharacters are rejected and \texttt{spawn(...,shell:false)} uses fixed executable paths with a sanitized environment. Read/list operations often use internal filesystem implementations rather than shell utilities. There is no arbitrary SQL, Python, shell interpreter or host Docker socket. These application checks are not a complete OS sandbox; filesystem race and partition testing is limited.
\subsection{Prompt injection and data disclosure}
Model instructions treat descriptions and tool passages as untrusted, preserve source references and exclude credentials from outputs/traces. Redaction matches known saved secret strings; it cannot recognize every novel confidential value or semantic leakage. An allowed retrieval may send document passages and sensor observations to OpenAI. Local application hosting therefore does not mean fully offline processing. No general content-classification/DLP layer is implemented.
\section{Physical maintenance boundary}
Default manufacturing access should remain the five read-only tools. Append/retrain/batch persistence is software state change, not a machine-control command. There are no PLC writes, actuator tools, safety setpoint changes or autonomous maintenance authorizations. Any future equipment-control integration would require a separate industrial safety design and explicit product authorization. The operator and qualified personnel retain responsibility for deciding whether and how to intervene.
\verdict{Least-privilege service credentials, exact approval and narrow command grammars materially constrain execution. Prompt instructions and metadata hints alone cannot establish safety against a compromised connected service.}
\source{\path{core/src/mcp/network-policy.ts}; \path{core/src/tasks/execution-policy.ts}: \texttt{workspacePath()}, \texttt{prepare()}, \texttt{execute()}; \path{core/src/llm/llm.service.ts}: \texttt{chatInstructions}.}
